\documentclass[10pt,a4paper]{amsart}
\usepackage[margin=19mm]{geometry}
\usepackage{amsmath,amssymb,mathtools}
\usepackage[scr=rsfs,cal=euler]{mathalfa}
\usepackage{xfrac}
\usepackage[unicode,hidelinks,bookmarks=false]{hyperref}
\newcommand{\ie}{i.e.\ }
\newcommand{\cf}{cf.\ }
\newcommand{\resp}{respectively\ }
\newcommand{\Subsection}{\subsection}
\providecommand{\nobreakdash}{\nobreak\hbox{-}}
\setlength{\emergencystretch}{2em}
\multlinegap0pt
\usepackage{bm}
\usepackage{bbm}
\usepackage{dsfont}
\usepackage{xspace}
\usepackage{cancel}
\usepackage{mathtools}
\usepackage{amsthm}
\makeatletter
\@ifundefined{theorem}{\newtheorem{theorem}{Theorem}}{}
\@ifundefined{lemma}{\newtheorem{lemma}[theorem]{Lemma}}{}
\@ifundefined{proposition}{\newtheorem{proposition}[theorem]{Proposition}}{}
\makeatother
\title[Evasive sets, twisted varieties, and container-clique trees]{Evasive sets, twisted varieties, and container-clique trees}
\author{Jeck Lim, Jiaxi Nie, Ji Zeng}
\date{Source: arXiv:2507.07594v1 (CC BY 4.0)}
\begin{document}
\maketitle

\noindent\emph{Source and licence.} Evasive sets, twisted varieties, and container-clique trees. Version arXiv:2507.07594v1 (CC BY 4.0), distributed under the Creative Commons Attribution 4.0 International licence, as recorded in the arXiv licence metadata for this exact version and shown on the abstract page https://arxiv.org/abs/2507.07594v1. Full rights notice: licenses/arxiv-2507.07594-CC-BY-4.0.txt.\par\smallskip

\noindent\textit{Editorial scope.} Counted unit: the Lemma whose source label is \texttt{counting\_supersaturation} in the article (source0.tex lines 262--266, source label \texttt{counting\_supersaturation}), delivered together with the article's own complete original proof of it (source0.tex lines 268--285, one \texttt{proof} environment of 18 lines). No mathematics is written, repaired or re-derived: statement and proof are the article's own text line for line. Required context retained uncounted: counting (lines 58--60); no-three-in-line\_key (lines 196--198); HCL (lines 205--214). Every definition, display and label of those lines is retained unchanged. Omitted from this package and not redistributed: 1--57, 61--195, 199--204, 215--261, 267--267, 286--356. Nothing inside a retained excerpt is omitted or altered: every retained body is delivered exactly as the article's lines are written, so no omission receipt of any kind accompanies this package. Citations: the retained excerpts carry no citation command at all, so no bibliography entry of the article is transcribed and no reference is redistributed. Reference adaptation: none was needed and none was made; every reference inside the delivered unit resolves inside it, so no unresolved reference or citation is produced. Printed slips kept literal: no printed word, symbol, hypothesis or number of the counted statement or of its proof was altered. Layout and class: the article's own class is \texttt{article} with packages including amssymb, amsmath, amsthm, hyperref, geometry; the wrapper is the lane's pinned \texttt{amsart} editorial wrapper at 10pt on A4, which restates the theorem-like environments the retained text needs and adds the article's own macro definitions that the retained text needs (none), with extra wrapper packages (bm, bbm, dsfont, xspace, cancel, mathtools). The delivered numbering is the unit's own. Assets: no retained span contains a figure environment, a graphics-inclusion command, a PDF-page inclusion, a table or a TikZ picture; no image file is copied into this package, the package declares no asset dependency and no image file is redistributed with it. Licence: version arXiv:2507.07594v1 of this article is distributed under the Creative Commons Attribution 4.0 International licence, as recorded in the arXiv licence metadata for this exact version and shown on the abstract page https://arxiv.org/abs/2507.07594v1. A full rights notice is delivered at licenses/arxiv-2507.07594-CC-BY-4.0.txt. Characters: every retained excerpt is pure ASCII, so the delivered engine, XeLaTeX with its default Latin Modern text font, reports no missing character.
\begin{theorem}\label{counting}
    There are at most $2^{O(q^{n-k})}$ many $(k,r)$-evasive sets in $\mathbb{F}^n_q$.
\end{theorem}

\begin{proposition}\label{no-three-in-line_key}
For $p \geq \log^3 q /\sqrt{q}$, we have $\alpha(\mathbb{F}^2_q,p) \leq (1+o(1))pq$ with high probability.
\end{proposition}

\begin{lemma}\label{HCL}
Let $r \ge 2$ be an integer and $c>0$ be sufficiently small with respect to $r$. If $\mathcal{H}=(V,E)$ is an $r$-uniform hypergraph and $0<\tau<1/2$ is a real number such that \begin{equation*}
        \Delta_i(\mathcal{H}) \leq c \cdot \tau^{i-1} \frac{|E|}{|V|} \quad\text{for all $2\leq i\leq r$,}
\end{equation*} then there exists a family $\mathfrak{C}$ of vertex subsets of $\mathcal{H}$ with the following properties:
\begin{itemize}
            \item[(a)] Every independent set of $\mathcal{H}$ is contained in some $C \in \mathfrak{C}$.
            \item[(b)] $\log |\mathfrak{C}| \leq c^{-1} \cdot \tau |V| \cdot \log(1/\tau)$.
            \item[(c)] For every $C \in \mathfrak{C}$, we have $|E(\mathcal{H}[C])| \leq (1 - c)|E|$.
\end{itemize}
\end{lemma}

\begin{lemma}\label{counting_supersaturation}
    For positive integers $k < n, r$ and real number $c > 0$, there exists a sufficiently large $\theta$ such that the following holds. If $P \subset \mathbb{F}_q^{n}$ satisfies $ |P| \geq \theta q^{n-k}$ and no $k$-flat intersects $P$ in more than $2|P|/\sqrt{q}$ points, then there exists a hypergraph $\mathcal{H}$ with $V(\mathcal{H}) = P$ and $E(\mathcal{H})$ consisting of some $(k,r)$-sets such that, for all $2\leq i\leq r$, \begin{equation*}
        \Delta_i(\mathcal{H}) \leq c \cdot \tau^{i-1} \frac{|E(\mathcal{H})|}{|V(\mathcal{H})|} \quad\text{with}\quad \tau = \frac{\theta q^{n-k}}{|P|q^\epsilon} ~\text{and}~ \epsilon = \frac{1}{2r},
    \end{equation*} and $\Delta_1(\mathcal{H})\leq \theta |E(\mathcal{H})|/|V(\mathcal{H})|$.
\end{lemma}

\begin{proof}[Proof of Theorem~\ref{counting}]

    We assume $n,r > k$ since the other cases are trivial. Let $\mathcal{H}$ be the hypergraph with $V(\mathcal{H}) = \mathbb{F}^n_q$ and $E(\mathcal{H})$ consisting of all $(k,r)$-sets. We build a container-clique tree $T$ for $\mathcal{H}$ through the following process. We start with $T$ having one root node labelled by $\mathbb{F}_q^n$; iteratively, if there is a node $x \in T$ with $|C^x_0| \geq 2\theta q^{n-k}$ for the $\theta$ in Lemma~\ref{counting_supersaturation} (with $c$ in Lemma~\ref{HCL}), we perform an operation on $x$ as described below; the process stops when no such node exists.

    Now, we describe the operation on a node $x$. We take $C = C^x_0$ first. If there exists a \textit{rich} $k$-flat $F$ in $\mathbb{F}_q^n$, that means, $|C \cap F| \geq |C^x_0|/\sqrt{q}$, then we delete the subset $C \cap F$ from $C$. We repeat this deletion step until either (i) $|C| < |C^x_0|/2$ or (ii) no rich $k$-flat exists. We use $K_1,K_2,\dots,K_\ell$ to denote the deleted subsets. Note that every $K_i$ is a clique of $\mathcal{H}$. If the deletion steps ended with (i), then we create a child $y$ for $x$ such that $y$ copies the label of $x$, but with $C^x_0$ replaced by $C$ and all $K_i$'s appended.

    If the deletion steps ended with (ii), we consider an auxiliary hypergraph $\mathcal{H}'$ output by applying Lemma~\ref{counting_supersaturation} to $C$ as $P$ and $c$ as in Lemma~\ref{HCL}. Here, $|C| \geq |C^x_0|/2 \geq \theta q^{n-k}$ and the absence of rich flats are used. Thereafter, we apply Lemma~\ref{HCL} to $\mathcal{H}'$ with $\tau$ as given in Lemma~\ref{counting_supersaturation} to obtain a family $\mathfrak{C}$ of subsets of $C$. For each $C' \in \mathfrak{C}$, we create a child $y$ for $x$ with $C^y_0 = C'$, $C^y_1 = C^x_1, \dots, C^y_{\ell_x}=C^x_{\ell_x}$, $C^y_{\ell_x+1} = K_1,\dots, C^y_{\ell_x+\ell} = K_\ell$. Thus, the node operation is described. This process will give us a container-clique tree $T$ by the same argument as in Proposition~\ref{no-three-in-line_key}.

    We argue that each node operation shrinks the size of $C^x_0$ by a constant factor. This is self-evident if the deletion steps of $x$ ended with (i). If the deletion steps of $x$ ended with (ii), Lemma~\ref{HCL}(c) tells us that $|E(\mathcal{H}'[C'])| \leq (1-c)|E(\mathcal{H}')|$ for each $C' \in \mathfrak{C}$ in the node operation of $x$. It is also guaranteed that $\Delta_1(\mathcal{H}')\leq \theta |E(\mathcal{H}')|/|V(\mathcal{H}')|$. Thus, we can compute \begin{equation*}
        (|C|-|C'|) \cdot \theta\frac{|E(\mathcal{H'})|}{|V(\mathcal{H'})|}\geq(|C|-|C'|)\cdot\Delta_1(\mathcal{H}')\geq|E(\mathcal{H'}[C\setminus C'])|\geq c|E(\mathcal{H}')|.
    \end{equation*} Simplifying the inequality above gives $|C'| \leq (1-c/\theta)|C|$, which means $|C^y_0| \leq (1-c/\theta)|C^x_0|$ for every child $y$ of $x$.

    Next, we examine some quantitative properties of $T$. Firstly, since each node operation shrinks the size of $C^x_0$ by a constant factor, the height of $T$ is at most $O(\log q)$. Secondly, since each node operation extends the label by at most $\sqrt{q}$ additional cliques, every node has label length at most $O(\sqrt{q}\log q)$. By our choice of $\tau$ and Lemma~\ref{HCL}(b), the degree of a node will be at most $2^{O(q^{n-k-\epsilon}\log q)}$ with $\epsilon=1/(2r)$ as in Lemma~\ref{counting_supersaturation}, which implies that $T$ has at most $2^{O(q^{n-k-\epsilon}\log^2 q)}$ many leaves. Lastly, we have $|C^x_0| < 2\theta q^{n-k}$ for all leaf $x$.

    Finally, we upper bound the number of $(k,r)$-evasive sets in $\mathbb{F}_q^n$, which is exactly the number $\aleph$ of independent sets in $\mathcal{H}$. We denote the number of leaves as $\nu$, the maximum size of $C^x_0$ for leaf $x$ as $\chi$, the maximum size of $C^x_i$ for $i\neq 0$ as $\kappa$, and the maximum label length $\ell_x$ as $\lambda$. Since an independent set and a clique intersect in less than $r$ vertices, we have \begin{equation*}
        \aleph \leq \nu \cdot \binom{\kappa}{r}^\lambda \cdot 2^{\chi+r\lambda}  \leq 2^{O(q^{n-k-\epsilon}\log^2 q)} \cdot \binom{q^2}{r}^{\sqrt{q}\log q} \cdot 2^{O(q^{n-k})} \leq 2^{O(q^{n-k})}.
    \end{equation*} Therefore, we conclude the proof.
\end{proof}

\end{document}
